\chapter{Termo de Abertura do Projeto}
\label{termo}

\section{Dados do projeto}
\begin{description}
    \item [Nome do Projeto:] Micromouse
    \item [Data de abertura:] 05/05/2026
    \item [Código:] 1-A
    \item [Patrocinador:] Universidade de Brasília
    % \item [Responsável:] ???
    \item [Gerente:] Ian Costa Guimarães/222014859/iancostag@gmail.com/6199299-8915
\end{description}

\section{Objetivos}
O objetivo central é projetar, construir e programar um robô móvel autônomo capaz de mapear e resolver labirintos de 4x4, 8x8 e 16x16 células, integrando-o a uma plataforma web para telemtria. Seguindo a regra SMART:
\begin{description}
    \item [\textit{Specific} (específico):] Desenvolver um robô físico (Micromouse) que resolva 3 labirintos distintos de forma autônoma e um sistema web para visualização de telemetria (trajeto, velocidade e bateria) em tempo real.
    \item [\textit{Measurable} (mensurável):] O sucesso será medido pela conclusão dos labirintos de até 16x16 células e pela recepção de pacotes de dados via UDP/WebSockets com latência inferior a 50ms.
    \item [\textit{Agreed} (acordado):] O projeto foi validado pelos líderes das quatro áreas (Software, Eletrônica, Energia e Estruturas) e pelo gerente geral, garantindo a integração dos subsistemas.
    \item [\textit{Realistic} (realista):] O projeto utiliza tecnologias acessíveis e consolidadas (ESP32, Motores N20, Sensores ToF) compatíveis com o orçamento e o laboratório de eletrônica da FGA.
    \item [\textit{Time Bound} (limitado no tempo):] A entrega final do protótipo funcional e do relatório técnico completo (ED1) está definida para 11 de maio de 2026.
\end{description}

\section{Mercado-alvo}

O projeto destina-se a instituições acadêmicas e de pesquisa em robótica móvel, além de competições tecnológicas promovidas pelo IEEE. Internamente, atende aos requisitos pedagógicos da disciplina de Projeto Integrador de Engenharia 1 da Faculdade de Ciências e Tecnologias em Engenharia (FCTE-UnB), servindo como plataforma de estudo para estudantes de Engenharia de Software, Eletrônica, Energia, Aeroespacial e Automotiva, promovendo entendimento holístico de um projeto de engenharia muldisciplinar real.

\section{Requisitos}

Os requisitos essenciais para a viabilidade do produto incluem:
\begin{itemize}
    \item \textbf{Dimensões Operacionais:} O robô deve possuir largura entre 7cm e 10cm para navegar livremente em corredores de 16,8cm.
    \item \textbf{Processamento:} Utilização de microcontrolador dual-core (ESP32) para processamento paralelo de controle e telemetria.
    \item \textbf{Navegação:} Sensores de distância do tipo Time-of-Flight (VL53L0X) para mapeamento preciso e independente da cor das superfícies.
    \item \textbf{Monitoramento:} Dashboard web responsivo que exiba a posição do robô em tempo real via Wi-Fi.
    \item \textbf{Autonomia:} Sistema de energia capaz de alimentar motores e comunicação sem fio durante o ciclo completo de resolução de 3 labirintos.
\end{itemize}

\section{Justificativa}

 O desenvolvimento do projeto Micromouse é justificado pela necessidade de integrar, em um único sistema, conhecimentos multidisciplinares das diferentes engenharias da Faculdade de Ciências e Tecnologias em Engenharia (FCTE-UnB). O projeto permite aplicar conceitos de sistemas embarcados, eletrônica, controle, programação, mecânica, banco de dados e desenvolvimento web em um contexto prático e competitivo, aproximando os estudantes de problemas reais de engenharia.

Além do caráter acadêmico, o projeto atende à crescente demanda por soluções autônomas capazes de realizar navegação inteligente em ambientes desconhecidos. Tecnologias semelhantes às utilizadas em competições Micromouse também estão presentes em aplicações industriais, robótica móvel, logística automatizada e veículos autônomos, tornando o projeto relevante como ferramenta de capacitação técnica e desenvolvimento profissional.

Outro fator que justifica o projeto é a oportunidade de desenvolver uma plataforma de monitoramento em tempo real integrada ao robô, que permite acompanhar métricas de desempenho, telemetria e comportamento do sistema durante a execução das tarefas. Essa funcionalidade amplia o potencial educacional do projeto, possibilitando análises posteriores, validação experimental e apoio ao processo de tomada de decisão durante os testes.

No contexto da disciplina de Projeto Integrador de Engenharia 1, o Micromouse também contribui para o desenvolvimento de competências relacionadas à gestão de projetos, trabalho em equipe, rastreabilidade de atividades, documentação técnica e integração entre áreas, competências fundamentais para a formação de engenheiros.

\section{Indicadores}

Os indicadores definidos para o projeto têm como objetivo medir o desempenho técnico do sistema desenvolvido, bem como avaliar sua aderência aos requisitos estabelecidos para a competição e para a disciplina. Os principais indicadores considerados são:

\begin{itemize}
    \item Taxa de sucesso na resolução dos labirintos de 4x4, 8x8 e 16x16 células;
    
    \item Tempo médio necessário para completar o percurso até a região objetivo;
    
    \item Latência média da transmissão de telemetria entre o robô e o dashboard web;
    
    \item Precisão do mapeamento do labirinto realizado pelo algoritmo de navegação;
    
    \item Número de colisões ou falhas de navegação durante os testes;
    
    \item Autonomia energética do sistema durante as execuções consecutivas;
    
    \item Taxa de perda de pacotes de telemetria durante a comunicação Wi-Fi;
    
    \item Velocidade média de deslocamento do robô durante a exploração;
    
    \item Quantidade de requisitos funcionais implementados e validados nos testes;
    
    \item Percentual de atividades concluídas dentro do prazo planejado no GitHub Projects.
\end{itemize}
% \section{Aprovações}

% \begin{tabular}{ l l }
%   \textbf{Patrocinador:} & \_\_\_\_\_\_\_\_\_\_\_\_\_\_\_\_\_\_\_\_\_\_\_\_\_\_\_\_\_\_\_\_\_\_ \\
%   & \\
%   \textbf{CEO:} & \_\_\_\_\_\_\_\_\_\_\_\_\_\_\_\_\_\_\_\_\_\_\_\_\_\_\_\_\_\_\_\_\_\_ \\
%   & \\
%   \textbf{Gerente:} & \_\_\_\_\_\_\_\_\_\_\_\_\_\_\_\_\_\_\_\_\_\_\_\_\_\_\_\_\_\_\_\_\_\_ \\
% \end{tabular}
